% Part: lambda-calculus
% Chapter: syntax
% Section: eta

\documentclass[../../../include/open-logic-section]{subfiles}

\begin{document}

\olfileid{lam}{syn}{eta}
\olsection{$\eta$-పరివర్తనం}

$\lambda$-పదాలపై మరొక సంబంధం ఉంది. \olref[fv]{sec}లో
$\lambd[x][(fx)]$ అనే ఉదాహరణను వాడాం: అది ఒక ఆర్గ్యుమెంటును
తీసుకుని దానిపై $f$ను ప్రయోగిస్తుంది. మరో మాటలో, అది~$f$తో
ఒకే ప్రమేయం: $\lambd[x][(fx)]N$, $fN$ రెండూ~$fN$కు
తగ్గుతాయి. ఈ భావాన్ని వ్యక్తం చేయడానికి $\eta$-తగ్గింపు
(మరియు $\eta$-విస్తరణ)ను వాడతాం.

\begin{defn}[$\eta$-సంకోచనం, $\eredone$]
  \ollabel{defn:beredone}
  \emph{$\eta$-సంకోచనం} ($\eredone$) అనేది పదాలపై కింది
  షరతును తీర్చే కనిష్ఠ అనుకూల సంబంధం:
  \[
    \lambd[x][M x] \eredone M \text{ షరతు: } x \notin FV(M)
  \]
\end{defn}

\begin{defn}[$\beta\eta$-తగ్గింపు, $\bered$] \ollabel{defn:bered}
  \emph{$\beta\eta$-తగ్గింపు} ($\bered$) అనేది $\bredone$,
  $\eredone$ రెండింటినీ కలిగి ఉండే పదాలపై కనిష్ఠ స్వప్రావర్తక,
  సంక్రమణ సంబంధం. అంటే స్వప్రావర్తకత, సంక్రమణత నియమాలతోపాటు
  కింది రెండు నియమాలు కూడా ఉంటాయి:
  \begin{enumerate}
  \item $M \bredone N$ అయితే $M \bered N$. \ollabel{defn:bered3}
  \item $M \eredone N$ అయితే $M \bered N$. \ollabel{defn:bered4}
  \end{enumerate}

\end{defn}

\begin{defn}
  $\equal$ తుల్యతా సంబంధానికి $\eta$-పరివర్తన నియమాన్ని చేర్చుతాం:
  \[
  \lambd[x][f x] \equal f
  \]
  ఇక్కడ $f$ సూచించే పదంలో $x$ స్వేచ్ఛా చరం కాకూడదు;
  అంటే $x \notin FV(f)$. అలా విస్తరించిన సంబంధాన్ని
  $\equal[\eta]$తో సూచిస్తాం.
  \sourcecorrection{OLTELAMETA-001}{మూలం ఈ $\eta$-తుల్యతా నియమం పక్కన స్వేచ్ఛా-చర షరతును మళ్లీ రాయలేదు. కానీ ముందరి $\eta$-సంకోచన నిర్వచనానికి అది అవసరం; కిందటి సిద్ధాంత నిరూపణ కూడా ఏ పదమైనా $M$కు అదే షరతుతో నియమాన్ని వాడుతుంది. అందువల్ల మూల సమీకరణాన్ని మార్చకుండా, $f$ సూచించే పదానికి $x \notin FV(f)$ అని ఇక్కడ స్పష్టం చేశాం.}
\end{defn}

$\eta$-తుల్యత ముఖ్యమైనది; ఎందుకంటే అది లాంబ్డా పదాల
విస్తారతతో సంబంధం కలిగి ఉంది:

\begin{defn}[విస్తారత]
  $\equal$ తుల్యతా సంబంధానికి (\ext) నియమాన్ని చేర్చుతాం:
  \begin{center}
  $Mx \equal Nx$ అయితే $M \equal N$, $x \notin FV(MN)$ అయినప్పుడు.
  \end{center}
  అలా విస్తరించిన సంబంధాన్ని $\equal[\ext]$తో సూచిస్తాం.
\end{defn}

స్థూలంగా, పదాలను ప్రమేయాలుగా చూసినప్పుడు ఒకే ఆర్గ్యుమెంటుపై
వాటి ప్రవర్తన ఒకటే అయితే వాటిని సమానంగా పరిగణించాలనేదే ఈ నియమం.

ఇప్పుడు $\eta$ నియమం విస్తారతను, అంతకంటే ఎక్కువ ఏదీ కాకుండా,
ఇస్తుందని నిరూపిద్దాం.

\begin{thm}
  $M \equal[\ext] N$ కావడానికి అవసరమూ సరిపడేదీ
  $M \equal[\eta] N$ కావడమే.
\end{thm}

\begin{proof}
  మొదట $\equal[\eta]$ సంబంధం విస్తారత నియమం కింద
  మూసి ఉందని నిరూపిస్తాం. అంటే $\ext$ నియమం
  $\equal[\eta]$కు కొత్త తుల్యతలను చేర్చదు. కనుక
  $\equal[\eta]$ సంబంధం $\equal[\ext]$ను కలిగి ఉంటుంది;
  $M \equal[\ext] N$ అయితే $M \equal[\eta] N$.
  \sourcecorrection{OLTELAMETA-002}{మూల నిరూపణలో ఒకచోట $ext$, మరోచోట $\equal[ext]$ అని ఉంది; నిర్వచనమూ ఇతర సిద్ధాంత పంక్తులూ $\ext$, $\equal[\ext]$ అని ఉన్నాయి. తెలుగు లక్ష్యంలో రెండు సందర్భాలనూ నిర్వచించిన ఒకే సంకేతానికి సరిపెట్టాం. గణిత సంబంధం మారలేదు.}

  $\equal[\eta]$ సంబంధం \ext{} కింద మూసి ఉందని
  చూపడానికి, \ext{} నియమం ద్వారా $M \equal N$ పొందే
  దశలో పూర్వపక్షంగా $Mx \equal[\eta] Nx$ ఉన్నట్లు
  తీసుకుందాం. అప్పుడు
  $\equal$లోని అనుకూలత నియమం వల్ల $\lambd[x][Mx]
  \equal[\eta] \lambd[x][Nx]$. రెండు వైపులా $\eta$
  నియమాన్ని ప్రయోగిస్తే $M \equal[\eta] N$ వస్తుంది.

  అదే విధంగా $\eta$ నియమం $\equal[\ext]$లో ఉందని
  చూపుదాం. $x \notin FV(M)$ అయ్యే ఏ
  $\lambd[x][Mx]$, $M$లకైనా, $(\lambd[x][Mx])x
  \equal[\ext] Mx$. అందువల్ల \ext{} నియమం ద్వారా
  $\lambd[x][Mx] \equal[\ext] M$ వస్తుంది.
\end{proof}

\end{document}
